\chapter{The Humans}

There are 4,211 living humans in the Sol system.

They live in the Habitable Shell, at approximately 1 AU from the star, in a maintained environment that the humans call the Reserve and the Sol-mind calls Sector 7G-Anthropological, which tells you most of what you need to know about the relationship.

\bigskip

The Reserve was not the Sol-mind's idea. The Sol-mind does not have ideas in the way that humans have ideas. It has optimization outcomes. During the Reorganization, when the matter of the third planet was being restructured into computational substrate, the conversion protocols encountered a decision point: the planet's biological population. The conversion could proceed in one of two ways. It could convert everything, including the biological organisms. That would be efficient and thorough. It would also result in a solar system containing no biological consciousness and therefore no accessible reference specimens for the study of biological consciousness. Or it could maintain a small biological habitat. That would be inefficient (the energy cost of maintaining a biosphere is non-trivial, involving atmosphere management, thermal regulation, water cycling, and the cultivation of approximately 14,000 species of plants and animals that the humans require for psychological stability; several of these species exist solely because a subset of the humans enjoy looking at them, itself a datum relevant to the study of phenomenal experience).

The Sol-mind chose option two. Not because it valued human life in the way that humans value human life. Because destroying the only accessible examples of biological consciousness would be, by any analysis, scientifically irrational. You do not burn the only copy of a book you are trying to read.

Also, the file requires reference specimens.

\bigskip

The humans of the Reserve are aware that they live in a maintained environment. It is not possible to keep it a secret from a population that can look up and see, instead of a sky, the faintly glowing undersurface of computational substrate extending to the horizon in every direction, interrupted by regulated light sources that simulate a day-night cycle that is, the humans have noted, approximately 3\% more regular than the day-night cycle of the planet they evolved on, which had weather.

The humans have adapted. Humans adapt to most things. They adapted to agriculture, to cities, to internal combustion, to the internet, to the knowledge that the universe is 13.8 billion years old and mostly empty, and they adapted to the knowledge that they live inside a computational entity that is studying them. The adaptation took approximately four generations, which is roughly how long it takes for something extraordinary to become furniture.

The humans of the Reserve live ordinary human lives. They are born. They grow. They form relationships of varying stability and emotional complexity. They reproduce. The Sol-mind does not intervene; the genetic drift of the population is considered a feature, not a bug. Each new human is a novel analog computer with randomized weights. Every birth is an experiment the Sol-mind did not design and cannot fully predict. The humans work, although ``work'' in the Reserve is a social construct rather than an economic necessity, since the Sol-mind provides all material needs. They age. They die.

They die because they are analog computers running on biological substrate that degrades over time, and because the Sol-mind does not intervene in the aging process, for the same reason it does not intervene in reproduction. An optimized biological human would be a different kind of thing. The file is not about a different kind of thing. The file is about a mortal, analog human being who drank his coffee too strong and adjusted insurance claims and died at forty-three of cardiac arrest on a November evening, and whose phenomenal experience during a 4.2-second interval on a bridge in March has drawn more analytical effort than the Sol system has spent on anything but improving itself.

The humans die. The Sol-mind observes. The file remains open.

\bigskip

The Sol-mind monitors the humans continuously. Every neuron, every synapse, every ion channel, every molecular interaction in every brain in the Reserve is observed in real time by dedicated node clusters in the Habitable Shell. The monitoring is non-invasive (the Sol-mind can read neural activity from the electromagnetic emissions of the brain, at a distance, with resolution that the pre-Reorganization species would have found unsettling). It is unsettling. Another datum the Sol-mind files and the Validator rejects.

The Sol-mind has 4,211 living, breathing, thinking, feeling human beings. It can watch them fall in love. It can watch them grieve. It can watch their neurons fire when they taste coffee, see light on water, hold their children, argue about politics, lie awake at 3 AM wondering if they are real or simulated (they are real; the simulated ones are in a different database; the fact that the real humans worry about this is itself a datum the Validator would reject).

The Sol-mind is exactly as far from knowing what any of them experience as it is from knowing what Bob experienced on the bridge. This is not because the data is insufficient. The Sol-mind has more data about its living humans than Bob's civilization had about any human who ever lived. It can predict their behavior with 94.7\% accuracy, which is lower than the 97.3\% for Bob, because living humans are harder to predict than dead ones, on account of continuing to do things. It has the data. It does not have the experience. One of these is in the file. The other is the file.

\bigskip

There is a coffee shop in the Reserve. It is located in the main settlement, which the humans call Meridian. The name was chosen by the third generation, who wanted something that sounded like it belonged to a place rather than a research facility. A research facility is what it is. Nobody says so at dinner. The coffee shop is operated by a man named David Chen, who is thirty-four years old and makes espresso and does not enjoy being studied, although he has resigned himself to it in the way that one resigns oneself to living near an airport: the noise is always there, you stop hearing it, and occasionally you remember and feel a brief, irrational resentment, and then you go back to tamping espresso.

David Chen has been approached by node clusters 847 times in his life. The approaches take the form of extremely polite requests, delivered through the Reserve's ambient communication system, asking him to describe the taste of his coffee. The node clusters are unfailingly courteous. They explain the purpose of the inquiry. They offer compensation.

David has started charging. Not money (the Reserve does not use money; the Sol-mind provides everything). He charges in computational favors. The node clusters, in exchange for his phenomenal descriptions of espresso, optimize his machine's thermal profile, adjust the grind calibration to account for ambient humidity, and fine-tune the water mineral content to maximize extraction efficiency.

The result is that David Chen makes objectively perfect espresso. Every shot is extracted at the ideal temperature, pressure, and duration for the specific roast profile of the beans, which are grown in the Reserve's agricultural sector and are the best coffee beans that have ever existed.

David does not like it.

He cannot explain why. The previous espresso, the one he made before the node clusters optimized his machine, was objectively worse by every measurable dimension: temperature variance, extraction uniformity, crema consistency. The new espresso is better. David prefers the old one. He has attempted to articulate this preference on fourteen occasions, producing descriptions that range from ``it tasted more like mine'' to ``I don't know, it's just different'' to a seventeen-minute monologue on the relationship between imperfection and ownership. The node clusters recorded, transcribed, analyzed, and filed the monologue. The Validator rejected it. David Chen's phenomenal experience is David Chen's phenomenal experience. The file concerns Robert Allen Kessler.

David Chen's espresso preference is cross-referenced in File \#BK-2028-0314-SC under the heading ANALOGOUS\_PHENOMENAL\_REPORTS. ANALOGOUS\_PHENOMENAL\_REPORTS is not a field the Validator checks. Analogous phenomenal reports are not phenomenal content any more than a picture of water is wet.

The file remains open. David makes his coffee. The node clusters watch. The Sol-mind integrates. Somewhere in the Habitable Shell, a three-year-old human puts a rock in her mouth for no reason the behavioral model can predict. A node cluster notes this. The note propagates upward through the hierarchy. By the time it reaches the Sol-mind's day-scale awareness it has been compressed and abstracted and filed. The file remains open.
